problem:  
Let \( E_{p,q} = SU(3)//S^1_{p,q} \) be the seven-dimensional Eschenburg space obtained by the free circle action
\[
z \cdot A = \operatorname{diag}(z^{p_1},z^{p_2},z^{p_3})\,A\,\operatorname{diag}(\overline{z}^{\,q_1},\overline{z}^{\,q_2},\overline{z}^{\,q_3}),
\]
where \(p,q\in\mathbb{Z}^3\) satisfy \(\sum_i p_i=\sum_i q_i\).  
Let \(L\) be a left-invariant metric on \(SU(3)\) that is right-invariant with respect to \(S^1_{p,q}\), and let \(g_{p,q}(L)\) be the descended metric on \(E_{p,q}\) obtained via the Riemannian submersion \(SU(3)\rightarrow E_{p,q}\).  
The Ricci flow \(\partial_t g=-2\,\mathrm{Ric}(g)\) starting from \(g(0)=g_{p,q}(L)\) preserves the residual \(T^2\)-symmetry of this metric.

**Research problem:**  
Does there exist an Eschenburg space \(E_{p,q}\) and an initial submersion metric \(g_{p,q}(L)\) with nonnegative sectional curvature such that the Ricci flow \(g(t)\) on \(E_{p,q}\) develops strictly positive sectional curvature at some finite time?  
Can Ricci flow thus serve as a mechanism that dynamically creates positivity of sectional curvature on an Eschenburg space, even when such positivity cannot be achieved by explicit smooth deformations of submersion metrics?

Why is it a "good" problem:  
This question is well posed and coherent with the geometry of Eschenburg spaces: the construction and symmetry description are accurate, and the invocation of \(T^2\)-symmetry relies on known isometric properties of Eschenburg metrics. It asks whether Ricci flow, an analytic process, can **generate positive curvature** in a geometric setting where such metrics are rare and often highly constrained.  

It sits at the intersection of **geometric analysis (Ricci flow evolution)**, **Riemannian submersion theory**, and **differential geometry of biquotients**, testing whether analytic flow phenomena can overcome algebraic obstructions in positive curvature constructions. A resolution would clarify the extent to which Ricci flow interacts with or transcends the algebraic structure of these biquotients, providing a new analytic lens on exotic examples that straddle the line between homogeneous and nonhomogeneous geometry.  

The statement is concise but far-reaching: “Can Ricci flow create positive curvature on an Eschenburg space?”—a question that is open, concrete, and analytically deep, embodying the principle of a “narrow entrance, profound depth” research problem.